% 功能与版式参考示例；自己的汇报由 agent 根据 outline.md 和本地素材生成 main.tex。
\documentclass[
      language=chinese,  % 主题文字：chinese / english
      cjk=auto,          % 中文支持：auto 随语言选择；也可设为 true / false
      bibstyle=ieee,     % 文献样式：ieee；例如可改为 authoryear
      bibsorting=none,   % 文献排序：none 按引用顺序；例如 nyt 按姓名、年份、标题
      mathfont=serif,    % 数学字体：serif / sans
      navigation=true,  % 页脚章节导航：true / false
      cornerwidth=4.8cm, % 角落引用默认宽度，填写长度
      layoutguides=false % 封面、章节页与致谢页的红线诊断：true / false
]{campusbeamer}

% ============================== 示例配置区 ==============================
% 上方列出全部模板选项及其默认值，直接修改即可；页面保持固定 16:9。

\campusbibresource{bibliography/refs.bib} % 只指定本次汇报的文献文件，无需自行加载 biblatex
% 版式微调时可打开此开关：封面、章节起始页和致谢页会叠加红色
% 定位框、对齐线、各部分字号及其实际 baselineskip；交付前保持关闭。
% 仅通过 documentclass 选项 layoutguides=true 开启；省略或设为 false 则关闭。

% !TeX root = ../example.tex
% 本次演示的标题、课题组与个人信息；直接在这里修改。
\title{Tsinghua Beamer 模板} % 主标题
\subtitle{完整功能示例} % 副标题
\course{MMLab@SIGS组会汇报} % 课程、课题组或会议名称
\author{范潇} % 汇报人
\IDnumber{xiaofan140@gmail.com} % 可选：学号或工号；取消注释后显示在姓名后
\advisor{王智\,副教授} % 可选：填写后在封面显示“指导教师：...”
\date{\today} % 汇报日期；也可填写固定日期
 % 标题、作者、导师与日期
% ============================ 示例配置区结束 ============================

\begin{document}
\maketitle

% 章节按这里的顺序组合；正文与演示片段统一放在 chapters/。
% !TeX root = ../example.tex
\section{基础页面与图像}
% 本示例每节显式插入起始页：图片直接在本次汇报中选择。
% 实际使用时，可按需省略目录页或起始页，也可交换两者顺序。
\sectioncoverpage[assets/sigs_building]{本节介绍页面布局、公式排版、图像适配与图注。}
\note{本节先说明正文区域和字号，再演示公式、代码与图片。章节背景是模板已有示例素材。}
\subsection{基础页面}
% !TeX root = ../example.tex
\begin{frame}{先确认大纲，再制作汇报}
      \begin{description}[最终交付]
            \item[准备内容] 在根目录 \texttt{outline.md} 写草稿，素材放入 \texttt{materials/}。
            \item[确认大纲] agent 整理章节、每页主旨、正文与备注分工，交由你审阅确认。
            \item[制作汇报] 确认后生成正文与备注，检查页面并导出 PDF/PPTX。
            \item[最终交付] \texttt{build/main.pdf} 和带逐页备注的 \texttt{build/main.pptx}。
      \end{description}
      \medskip
      对 agent 说：根据大纲制作 PPT。
\end{frame}
 % 推荐用法：思路与素材交给 agent

% 红色参考线只显示在这一页，不改变其他页面的主题设置。
% Diagnostic instrumentation is local to this frame; the theme is unchanged.
\begingroup
  % Horizontal dimensions use mm; vertical dimensions use the named font's
  % baseline grid (header: tiny/Large/normalsize; footer: scriptsize).
\newcommand{\layoutGuideDistance}[2]{%
  \begingroup
    \pgfextractx{\dimen0}{\pgfpointdiff{\pgfpointanchor{#1}{center}}{\pgfpointanchor{#2}{center}}}%
    \pgfextracty{\dimen2}{\pgfpointdiff{\pgfpointanchor{#1}{center}}{\pgfpointanchor{#2}{center}}}%
    \pgfmathparse{veclen(\the\dimen0,\the\dimen2)/(1mm)}%
    \pgfmathprintnumber[fixed,precision=2]{\pgfmathresult}\,mm%
  \endgroup
}
% Vertical dimensions are expressed in the same grid unit used by the
% navigation and footer rows. Horizontal dimensions continue to use mm.
\newcommand{\layoutGuideBaselineDistance}[3][\campusfooterbaseline]{%
  \begingroup
    \pgfextracty{\dimen2}{\pgfpointdiff{\pgfpointanchor{#2}{center}}{\pgfpointanchor{#3}{center}}}%
    \pgfmathparse{abs(\the\dimen2)/(1mm) / ((\the#1)/(1mm))}%
    \pgfmathprintnumber[fixed,precision=2]{\pgfmathresult}\,\mbox{baselineskip}%
  \endgroup
}
\newcommand{\layoutGuideHorizontalBaselineDistance}[3][\campusheadertextgap]{%
  \begingroup
    \pgfextractx{\dimen0}{\pgfpointdiff{\pgfpointanchor{#2}{center}}{\pgfpointanchor{#3}{center}}}%
    \pgfmathparse{abs(\the\dimen0)/(1mm) / ((\the#1)/(1mm))}%
    \pgfmathprintnumber[fixed,precision=2]{\pgfmathresult}\,\mbox{baselineskip}%
  \endgroup
}
\begin{frame}{版式对齐检查}
  % Draw last, above Beamer's headline/footline, on this page only.
  \AddToHookNext{shipout/foreground}{%
    \begin{tikzpicture}[remember picture,overlay]
      \tikzset{layout guide/.style={red,line width=.4pt},
        layout label/.style={text=red,font=\tiny,fill=white,inner sep=.7pt},
        layout dimension/.style={layout guide,<->}}
      % Both transparent assets are square and cropped to the emblem bounds.
      \coordinate (emblem-left) at
        ([xshift=\dimexpr\campuslogoleftskip+\campuslogopadding\relax,
          yshift=\dimexpr-\campuslogoheight+\campuslogopadding+.5\campuslogoimagewidth\relax]
          current page.north west);
      \coordinate (emblem-right) at ([xshift=\campuslogoimagewidth]emblem-left);
      \coordinate (emblem-bottom) at
        ([xshift=\dimexpr\campuslogoleftskip+.5\campuslogowidth\relax,
          yshift=\dimexpr-\campuslogoheight+\campuslogopadding\relax]
          current page.north west);
      \draw[layout guide]
        ([yshift=.5\campuslogoimagewidth]emblem-left) rectangle
        ([yshift=-.5\campuslogoimagewidth]emblem-right);
      \draw[layout dimension] ([xshift=-\campuslogopadding]emblem-left) -- (emblem-left);
      \draw[layout dimension] (emblem-right) -- ([xshift=\campuslogopadding]emblem-right);
      \draw[layout dimension] (emblem-bottom) -- ([yshift=-\campuslogopadding]emblem-bottom);
      \node[layout label,anchor=north west] at
        ([xshift=\campuslogoleftskip,yshift=\dimexpr-\campuslogoheight-2mm\relax]
          current page.north west)
        {校徽左右/下方留白：\the\campuslogopadding};
      \foreach \offset/\tag in {
        \campuslogoleftskip/V1,
        {\dimexpr\campuslogoleftskip+\campuslogowidth\relax}/V2,
        {\dimexpr\campuslogoleftskip+\campusframetitleleftskip\relax}/V3,
        {\dimexpr\paperwidth-\campuslogoleftskip\relax}/V4}{%
        \coordinate (\tag) at ([xshift=\offset]current page.north west);
        \draw[layout guide] (\tag) -- (\tag |- current page.south);
        \node[layout label,anchor=north] at
          ([yshift=-16mm]\tag) {\tag};
      }
      \draw[layout guide,dashed] (current page.north) -- (current page.south);
      \node[layout label,anchor=north] at ([yshift=-16mm]current page.north) {中轴};

      % H1, H2 and H3 are the actual boundaries of the three header grid
      % rows: after the empty row, after the title row, and after the
      % subtitle row.
      \coordinate (H1) at ([yshift=-\campusheaderemptybaseline]current page.north west);
      \coordinate (H2) at ([yshift=-\dimexpr\campusheaderemptybaseline+\campusheaderbaseline\relax]current page.north west);
      \coordinate (H3) at ([yshift=-\dimexpr\campusheaderemptybaseline+\campusheaderbaseline+\campusheadersubtitlebaseline\relax]current page.north west);
      \coordinate (H5) at ([yshift=\dimexpr\campusnavigationfooterheight+\campusfootlineheight+\campusfootlinedepth\relax]current page.south west);
      \coordinate (H6) at ([yshift=\dimexpr\campussectionnavigationheight+\campusfootlineheight+2\campusfootlinedepth\relax]current page.south west);
      \coordinate (H7) at ([yshift=\dimexpr\campusfootlineheight+\campusfootlinedepth\relax]current page.south west);
      \coordinate (H8) at (current page.south west);
      % H4 is the body safety boundary: keep a baseline-grid breathing room
      % above the first footer band (H5).
      \coordinate (H4) at ([yshift=2.4\campusfooterbaseline]H5);
      \foreach \tag in {H1,H2,H3,H4,H5,H6,H7}{%
        \draw[layout guide] (\tag) -- (current page.east |- \tag);
        \node[layout label,anchor=east] at
          ([xshift=-2mm]current page.east |- \tag) {\tag};
      }
      \coordinate (safe-a) at ([xshift=15mm]current page.west |- H4);
      \coordinate (safe-b) at (safe-a |- H5);
      \draw[layout dimension] (safe-a) -- (safe-b)
        node[midway,right,layout label] {\layoutGuideBaselineDistance{safe-a}{safe-b}};
      % An inset half-stroke keeps the bottom boundary visible in the PDF.
      \draw[layout guide] ([yshift=.2pt]current page.south west)
        -- ([yshift=.2pt]current page.south east);

      % Show all three header grid rows separately.  These coordinates are
      % layout boundaries, so glyph ascenders and descenders do not change the
      % reported one-baseline lengths.
      \coordinate (header-top) at (current page.north west);
      \foreach \upper/\lower/\pos/\label/\unit in {
        header-top/H1/.70/上方空行 (\string\tiny)/\campusheaderemptybaseline,
        H1/H2/.78/标题 (\string\Large)/\campusheaderbaseline,
        H2/H3/.86/副标题 (\string\normalsize)/\campusheadersubtitlebaseline}{%
        \coordinate (header-a) at ([xshift=\pos\paperwidth]current page.west |- \upper);
        \coordinate (header-b) at (header-a |- \lower);
        \draw[layout dimension] (header-a) -- (header-b)
          node[midway,left,layout label] {\label：\layoutGuideBaselineDistance[\unit]{header-a}{header-b}};
      }

      % Horizontal dimensions occupy the open area below the diagnostic text.
      \coordinate (logo-a) at ([yshift=-63mm]V1);
      \coordinate (logo-b) at (V2 |- logo-a);
      \draw[layout dimension] (logo-a) -- (logo-b)
        node[midway,above,anchor=south west,layout label] {校徽色块：%
          \layoutGuideHorizontalBaselineDistance[\campuslogowidth]{logo-a}{logo-b}
          (\string\Huge)};
      \coordinate (gap-a) at ([yshift=-67mm]V2);
      \coordinate (gap-b) at (V3 |- gap-a);
      \draw[layout dimension] (gap-a) -- (gap-b)
        node[midway,above,anchor=south west,layout label] {%
          \layoutGuideHorizontalBaselineDistance{gap-a}{gap-b} (\string\normalsize)};
      \coordinate (body-a) at ([yshift=-70mm]V1);
      \coordinate (body-b) at (V4 |- body-a);
      \draw[layout dimension] (body-a) -- (body-b)
        node[midway,above,layout label] {正文宽度：\layoutGuideDistance{body-a}{body-b}};
      \coordinate (margin-a) at ([yshift=-72mm]current page.north west);
      \coordinate (margin-b) at (V1 |- margin-a);
      \draw[layout dimension] (margin-a) -- (margin-b)
        node[midway,above,anchor=south west,layout label] {%
          \layoutGuideHorizontalBaselineDistance[\campusheaderbaseline]{margin-a}{margin-b} (\string\Large)};
      \coordinate (margin-c) at (V4 |- margin-a);
      \coordinate (margin-d) at (current page.east |- margin-a);
      \draw[layout dimension] (margin-c) -- (margin-d)
        node[midway,above,anchor=south east,layout label] {%
          \layoutGuideHorizontalBaselineDistance[\campusheaderbaseline]{margin-c}{margin-d} (\string\Large)};

      % Footer heights are read from the theme, so changes remain in sync.
      % Callouts above the bands keep their text clear of the navigation.
      \foreach \upper/\lower/\pos/\label in {
        H5/H6/.55/subsection (\string\scriptsize),
        H6/H7/.72/section (\string\scriptsize),
        H7/H8/.76/信息行 (\string\scriptsize)}{%
        \coordinate (band-a) at ([xshift=\pos\paperwidth]current page.west |- \upper);
        \coordinate (band-b) at (band-a |- \lower);
        \draw[layout dimension] (band-a) -- (band-b)
          node[midway,right,layout label] {\label：\layoutGuideBaselineDistance{band-a}{band-b}};
      }
    \end{tikzpicture}%
  }%
  \vspace*{11mm}
  \hspace*{-2pt}%
  \begin{minipage}{\dimexpr\textwidth-6pt\relax}
    \begin{description}[竖线]
      \item[竖线] V1 正文左界；V2 校徽色块右界；V3 页眉文字起点；V4 正文右界。
      \item[页眉] H1/H2/H3 为三行网格分界线；箭头以对应字号的 baselineskip 为单位。
      \item[页脚] H5--H7 为三行色带的上边界；间距以 \texttt{\string\scriptsize} 字号的 baselineskip 为单位。
      \item[参考] H4 为正文安全线，安全余量按页脚基线网格设置；虚线为页面中轴。
    \end{description}
  \end{minipage}
\end{frame}
\note{红线用于核对校徽、标题、正文和页脚的对齐。正式汇报通常不保留这张诊断页。}
\endgroup

% Read sizes from the active LaTeX commands rather than duplicating pt values.
\begingroup
\makeatletter
\newcommand{\fontGuideSize}[1]{%
  \begingroup #1\edef\fontGuideValue{\f@size}%
    \footnotesize\fontGuideValue\,pt\endgroup}
\makeatother
\newcommand{\fontGuideSample}[1]{%
  \makebox[24mm][l]{{\color{maincolor}#1 字 Aa}}}
\begin{frame}{模板字号与用途}
  \footnotesize
  \setlength{\tabcolsep}{0pt}
  \renewcommand{\arraystretch}{1.05}
  \begin{tabular}{@{}l@{\hspace{2mm}}p{28mm}@{\hspace{2mm}}p{15mm}@{\hspace{2mm}}p{72mm}@{}}
    \textbf{实际大小} & \textbf{字号命令} & \textbf{字号} & \textbf{主要用途} \\[.15mm]
    \fontGuideSample{\Huge\bfseries} & \texttt{\string\Huge} & \fontGuideSize{\Huge} &
      封面/章节/致谢页主标题 \\[.15mm]
    \fontGuideSample{\LARGE} & \texttt{\string\LARGE} & \fontGuideSize{\LARGE} &
      单行页眉（目录/参考文献等）；致谢页两行感谢语 \\[.15mm]
    \fontGuideSample{\Large\bfseries} & \texttt{\string\Large} & \fontGuideSize{\Large} &
      双行页眉上行、section 编号；封面／致谢页副标题 \\[.15mm]
    \fontGuideSample{\large} & \texttt{\string\large} & \fontGuideSize{\large} &
      可选的正文强调字号 \\[.15mm]
    \fontGuideSample{\normalsize} & \texttt{\string\normalsize} & \fontGuideSize{\normalsize} &
      正文/页眉下行；课题组/个人信息；小节导读 \\[.15mm]
    \fontGuideSample{\small} & \texttt{\string\small} & \fontGuideSize{\small} &
      提示框正文、二级列表；本例参考文献 \\[.15mm]
    \fontGuideSample{\footnotesize} & \texttt{\string\footnotesize} & \fontGuideSize{\footnotesize} &
      提示框标题、三级列表、图表题注 \\[.15mm]
    \fontGuideSample{\scriptsize} & \texttt{\string\scriptsize} & \fontGuideSize{\scriptsize} &
      三行页脚（两行导航与信息行） \\[.15mm]
    \fontGuideSample{\tiny} & \texttt{\string\tiny} & \fontGuideSize{\tiny} &
      红色标注、角落引用、论文元信息
  \end{tabular}
  \par\smallskip
  {\scriptsize 字号由当前文档实际读取，单位为 TeX pt。粗细由元素样式决定；数学上下标自动缩小。\par
  参考文献字号可通过命令参数调整；使用 shrink 时，整页缩放可能使最终显示更小。\par}
\end{frame}
\note{比较不同字号的实际显示。正文以正常字号为主，代码使用 small；内容过多时优先拆页。}
\endgroup


% section 组织整节内容，subsection 自动显示在页眉上方。
% frame 的标题参数即 frametitle，只填写本页标题，无需额外设置副标题。
\begin{frame}{三级内容结构}
      \begin{description}[Subsection]
            \item[Section] 组织整节内容，用于起始页和页脚导航。
            \item[Subsection] 组织节内主题，自动显示在页眉上方。
            \item[Frametitle] 概括本页内容，在页眉下方以小字号显示。
            \item[连续页面] 同一 subsection 下可放多页，每页只需填写自己的标题。
      \end{description}
\end{frame}
\note{说明 section 管整节，subsection 显示在页眉上行，frametitle 概括当前页。转入具体内容示例。}

% 公式与图像沿用本 subsection，保留单小节的导航示例。
% Formula examples use the amsmath environments already loaded by Beamer.
\begin{frame}{行内公式与多行推导}
  对样本 $\mathbf{x}_i\in\mathbb{R}^{d}$，用线性模型预测标量目标 $y_i$：
  \begin{align*}
    \widehat{y}_i &= \mathbf{w}^{\mathsf T}\mathbf{x}_i+b, \\
    \mathcal{L}(\mathbf{w},b)
      &= \frac{1}{N}\sum_{i=1}^{N}(\widehat{y}_i-y_i)^2
         +\lambda\lVert\mathbf{w}\rVert_2^2, \\
    \nabla_{\mathbf{w}}\mathcal{L}
      &= \frac{2}{N}\sum_{i=1}^{N}(\widehat{y}_i-y_i)\mathbf{x}_i
         +2\lambda\mathbf{w}.
  \end{align*}
  \begin{notebox}[排版要点]
    行内公式用 \texttt{\$...\$}；多行公式用 \texttt{align*}，
    在等号前放置 \texttt{\&} 对齐，星号表示不显示编号。
  \end{notebox}
\end{frame}
\note{沿等号解释预测、损失与梯度，强调多行推导应对齐关键运算符。这是数学排版示例。}

\begin{frame}{矩阵、分段函数与公式说明}
  \begin{columns}[T,onlytextwidth]
    \begin{column}{.47\textwidth}
      \textbf{矩阵：批量预测}
      \[
        \mathbf{X}=\begin{bmatrix}
          \mathbf{x}_1^{\mathsf T}\\
          \vdots\\
          \mathbf{x}_N^{\mathsf T}
        \end{bmatrix},\qquad
        \widehat{\mathbf{y}}=\mathbf{X}\mathbf{w}+b\mathbf{1}.
      \]
      \texttt{bmatrix} 自动生成方括号；矩阵中的列用 \texttt{\&} 分隔。
    \end{column}
    \begin{column}{.47\textwidth}
      \textbf{分段函数：ReLU}
      \[
        f(z)=\begin{cases}
          z, & z\geq 0,\\
          0, & z<0.
        \end{cases}
      \]
      \texttt{cases} 排列函数值与条件；公式中的文字可用 \texttt{\string\text\{...\}}。
    \end{column}
  \end{columns}
  \medskip
  \begin{tipbox}[符号与说明]
    $N$ 是样本数，$d$ 是特征维数；$\mathbf{1}\in\mathbb{R}^{N}$ 为全一向量。
    公式与符号解释放在同一页，便于阅读。
  \end{tipbox}
\end{frame}
\note{左侧检查矩阵间距，右侧检查分段条件。下方强调符号定义应在使用前说明。}


% !TeX root = ../example.tex
% Keep these examples within the existing subsection/navigation structure.
\begin{frame}[fragile]{代码展示：Python}
\begin{campuscode}[language=Python,numbers,caption={保留缩进，便于讲解函数结构。}]{normalize.py}
def normalize(values):
    # 将数值归一化，保留清晰的缩进
    total = sum(values)
    if total == 0:
        return [0.0 for value in values]
    return [value / total for value in values]

print(normalize([2, 3, 5]))
\end{campuscode}
\small 标题栏标识文件与语言；下方说明补充阅读提示。
\end{frame}

\begin{frame}[fragile]{代码展示：终端命令与输出}
\begin{campuscode}[language=bash,caption={静态展示命令与输出，不在编译时执行。}]{Terminal}
$ python -m venv .venv
$ source .venv/bin/activate
$ python normalize.py
[0.2, 0.3, 0.5]
\end{campuscode}
\small 用命令提示符区分输入与输出。
\end{frame}

\begin{frame}[fragile]{代码展示：从文件读取}
\campusinputcode[language=Python,description={从文件读取并高亮选定源码。}]
  {normalize.py}{chapters/code/normalize.py}
\small 用 \verb|\campusinputcode[language=Python]{文件名}{文件路径}| 复用源文件。
\end{frame}


\begin{frame}[fragile]{自动适配素材}
      \begin{columns}[T,onlytextwidth]
            \begin{column}{.48\textwidth}
                  \textbf{适配素材并添加图注}
                  \begin{verbatim}
\begin{figure}
  \fitgraphic[
    width=.42,
    height=.40
  ]{assets/sigs_wordmark.png}
  \caption{标志素材的适配效果}
\end{figure}
                  \end{verbatim}
                  宽、高分别是整页宽、高的比例，图片在范围内保持纵横比。
            \end{column}
            \begin{column}{.48\textwidth}
                  \begin{figure}
                        \fitgraphic[width=.42,height=.40]{assets/sigs_wordmark.png}
                        \caption{标志素材的适配效果}
                  \end{figure}
            \end{column}
      \end{columns}
\end{frame}
\note{对照左右两栏：左侧是插图源码，右侧是实际输出。宽高参数表示整页尺寸的比例，图片保持纵横比。}

\begin{frame}[fragile]{自动适配图片与图注}
      \begin{columns}[T,onlytextwidth]
            \begin{column}{.48\textwidth}
                  \textbf{同时设置图像与图注}
                  \begin{verbatim}
\fitfigure[
  width=.42,
  height=.28
]{\schoolwordmarkonlight}
 {横向标志素材的适配效果}
                  \end{verbatim}
                  最后一项填写图注，留空则不显示。
            \end{column}
            \begin{column}{.48\textwidth}
                  \fitfigure[width=.42,height=.28]{\schoolwordmarkonlight}%
                        {横向标志素材的适配效果}
            \end{column}
      \end{columns}
\end{frame}
\note{fitfigure 同时提供适配与图注。最后一个参数留空时不显示图注，便于简洁展示。}

\begin{rightimageframe}[image width=.40]{右侧满高图片}{assets/tsinghua_temple.png}
      \begin{description}[高度]
            \item[高度] 图片从页面顶端延伸到导航上边缘。
            \item[宽度] \texttt{image width} 控制右侧图片区域占页面宽度的比例。
            \item[正文] 左侧仍可放置列表、公式和提示框。
      \end{description}
\end{rightimageframe}
\note{检查右侧图片从页顶延伸至导航上沿，正文留在左侧。图片路径由当前汇报指定。}

\begin{frame}[fragile]{图像命令速查}
      \begin{description}[适配素材]
            \item[适配素材] \verb|\fitgraphic[width=.8,height=.6]{file}|
            \item[添加图注] \verb|\fitfigure[width=.8,height=.6]{file}{caption}|
            \item[右侧图片] \verb|\begin{rightimageframe}{标题}{file}|
            \item[PDF 页面] 可选参数 \texttt{page=2} 指定 PDF 的第 2 页。
      \end{description}
\end{frame}
\note{回顾三个常用图像入口；多页 PDF 可通过 page 参数选页。之后转入论文引用示例。}

% !TeX root = ../example.tex
\section{论文讲解与引用}
\sectiontocpage
\sectioncoverpage[assets/tsinghua_door]{本节演示论文讲解、文献字段与引用的显示规则。}
\subsection{论文讲解}

\begin{paperframe}{论文页面示例}{Zadeh2017TensorFN,Liu2018EfficientLM}[bottom-left][5.5cm]
      \begin{description}[引用]
            \item[正文] 放置论文要点、方法图或结果；每页聚焦一个主要信息。
            \item[页眉] 上行沿用当前小节，下行显示本页标题。
            \item[引用] 文献键关联 \texttt{bibliography/refs.bib}，本页自动在左下角列出两篇文献。
      \end{description}
      \medskip
      下一页展示生成本页的命令；文献仅用于演示引用排版。
\end{paperframe}

\begin{frame}[fragile]{论文页面示例：对应源码}
      \begin{verbatim}
\begin{paperframe}{论文页面示例}
  {Zadeh2017TensorFN,Liu2018EfficientLM}
  [bottom-left][5.5cm]
  % 论文要点、方法图或实验结果
\end{paperframe}
      \end{verbatim}
      \begin{description}[文献键]
            \item[页眉] 自动沿用当前 subsection；论文标题作为本页 frametitle。
            \item[文献键] 逗号分隔多篇文献，自动显示引用。
            \item[位置] 末尾可选 \texttt{top-right}（默认）、\texttt{bottom-left} 或 \texttt{bottom-right}。
            \item[宽度] 位置之后可继续填写引用宽度，例如 \texttt{[5.5cm]} 或 \texttt{[0.35\string\paperwidth]}。
      \end{description}
\end{frame}

\subsection{文献引用}

\begin{frame}{角落引用位置}
      \setcornercitewidth{5.5cm} % 仅本页引用加宽，下一页恢复全局宽度
      \begin{description}[右上]
            \item[右上] 默认位置支持用逗号分隔多篇文献。
            \item[左下] 最后一行引用与页脚上边缘保持设定距离。
            \item[右下] 适合左侧已有正文内容的页面。
      \end{description}

      \cornercite{Zadeh2017TensorFN,Liu2018EfficientLM}
      \cornercite[bottom-left]{Zadeh2017TensorFN}
      \cornercite[bottom-right]{Liu2018EfficientLM}
\end{frame}

\begin{frame}[fragile]{角落引用用法}
      \begin{verbatim}
\setcornercitewidth{5.5cm}
\cornercite{key1,key2}
\cornercite[bottom-left]{key1}
\cornercite[bottom-right]{key2}
      \end{verbatim}
      \begin{description}[引用宽度]
            \item[引用宽度] 调整命令放在引用之前，仅影响当前页。
            \item[文献信息] 在 \texttt{bibliography/refs.bib} 中维护，汇总页自动收集已引用条目。
      \end{description}
\end{frame}

% These rules describe the theme's current biblatex/IEEE configuration.
\begin{frame}{从文献字段到角落引用}
  \begin{description}[\texttt{shorttitle}]
    \item[\texttt{author}] 默认输出作者；姓名缩写和多人省略由引用样式决定。
    \item[\texttt{year}] 默认接在作者后，显示发表年份。
    \item[\texttt{title}] 接在作者、年份后，显示论文完整标题。
    \item[\texttt{shorttitle}] 仅在缺少 \texttt{title} 时作为替代标题。
  \end{description}
  \medskip
  本页右上角使用 Liu 条目，展示未设置 \texttt{usera} 时的默认输出。
  \cornercite{Liu2018EfficientLM}
\end{frame}
\note{区分标准作者、年份和标题字段；没有自定义短格式时，角落引用从标准字段生成。}

\begin{frame}[fragile]{自定义短文本：usera 的作用}
  以下摘自 \texttt{bibliography/refs.bib} 的 Zadeh 条目（其余字段省略）：
  \begin{verbatim}
usera      = {Zadeh et al., EMNLP 2017},
title      = {Tensor Fusion Network for
              Multimodal Sentiment Analysis},
shorttitle = {Tensor Fusion Network},
year       = {2017},
  \end{verbatim}
  \begin{itemize}
    \item 设置 \texttt{usera} 后：角落显示“自定义文本 + 完整标题”。
    \item 此时不会另加作者和年份；年份、会议简称需写在 \texttt{usera} 中。
    \item 本条目有 \texttt{title}，因此不会改用 \texttt{shorttitle}。
  \end{itemize}
  \cornercite{Zadeh2017TensorFN}
\end{frame}
\note{usera 自定义角落引用的短格式，不替换文末数字标签。不要把机构等未核实信息直接填入示例。}

\begin{frame}[fragile]{论文讲解页：补充论文信息}
  在文献条目中补充以下字段，\texttt{paperframe} 会自动追加信息：
  \begin{verbatim}
userb = {<论文署名机构>},
userc = {<会议或期刊，发表年月；预印本填写 arXiv>},
usere = {<引用次数>},
userf = {<统计来源，截至 YYYY-MM-DD>},
  \end{verbatim}
  \begin{itemize}
    \item 尖括号为填写提示；未核实的信息可以省略，对应行自动隐藏。
    \item 信息在同一段中连续显示，以竖线分隔，不额外占用多行。
    \item 普通角落引用与文末参考文献保持原格式。
  \end{itemize}
\end{frame}
\note{机构、发表状态、引用量和查询日期都是可选字段。此处尖括号内容只是格式占位符；引用数必须保留来源和日期。}

\begin{paperframe}[字段实际呈现]{论文信息字段示例}{Zadeh2017TensorFN}[bottom-left][7cm]
  \begin{itemize}
    \item 左下角沿用普通引用的灰色文字、字号和紧凑行距。
    \item 论文标题仍由引用中的 \texttt{title} 自动输出。
    \item 机构、发表 venue 或 arXiv、引用量和统计日期来自 \texttt{userb}、\texttt{userc}、\texttt{usere}、\texttt{userf}。
    \item 本页尖括号内容是待填写占位符，替换后即可用于正式汇报。
  \end{itemize}
\end{paperframe}
\note{查看附加字段的实际效果。此页用于排版演示，隐藏未填字段可以保持正式汇报简洁。}

\begin{frame}{文末参考文献：保留完整出处}
  \begin{description}[\texttt{booktitle}]
    \item[\texttt{author}] 按当前 IEEE 样式显示作者，名字使用首字母。
    \item[\texttt{title}] 显示论文题名；\texttt{usera} 不会替换这部分。
    \item[\texttt{booktitle}] 对 \texttt{@inproceedings} 条目，显示会议论文集名称。
    \item[\texttt{year}] 显示发表年份；页码、DOI 等按条目类型与样式输出。
  \end{description}
  \begin{notebox}[角落与文末的区别]
    角落引用不自动输出会议、页码或 DOI。
    文末页从已引用条目生成完整列表，按首次引用顺序编号；
    \texttt{usera} 不会改变这些数字标签。
  \end{notebox}
\end{frame}
\note{文末参考文献保留标准出处。角落短引用与文末完整条目承担不同作用。}

\begin{frame}[fragile]{文献键与编译流程}
  \begin{verbatim}
\addbibresource{bibliography/refs.bib}  % 导言区加载数据库
\cornercite{Zadeh2017TensorFN}
\referencespage[break][\small][onecolumn]
  \end{verbatim}
  \begin{description}[文献键]
    \item[文献键] \texttt{@inproceedings\{Zadeh2017TensorFN,...\}} 中的
      \texttt{Zadeh2017TensorFN} 是查找标识，不是打印出来的标题。
    \item[作者名] 多位作者在 \texttt{author} 字段内用 \texttt{and} 分隔。
    \item[编译] 使用 XeLaTeX 与 Biber；修改 \texttt{.bib} 后重新运行 \texttt{latexmk}。
  \end{description}
  \medskip
  只存入数据库但未引用的条目，默认不会出现在文末参考文献中。
\end{frame}
\note{按文献资源、正文引用、文末汇总的顺序回顾流程。biblatex 与 biber 由构建入口统一处理。}


% !TeX root = ../example.tex
\section{提示框}
\sectioncoverpage[assets/sigs_dom]{本节展示提示框样式与使用方法。} % 原图右边缘贴齐页面右边缘
\note{转入提示框：用有限的视觉强调区分建议、补充说明和警告。}
\subsection{样式演示}

\begin{frame}{提示、备注与警告}
      \begin{tipbox}[提示]
            使用 \texttt{tipbox} 放置操作建议，例如推荐的实验配置。
      \end{tipbox}

      \begin{notebox}[备注]
            使用 \texttt{notebox} 放置补充说明，例如数据来源或实验条件。
      \end{notebox}

      \begin{alertbox}[警告]
            使用 \texttt{alertbox} 强调易错点，例如复现时需要核对的设置。
      \end{alertbox}
\end{frame}
\note{比较提示、备注与警告的功能色。每个框只放一个需要强调的信息。}

% !TeX root = ../example.tex
\begin{frame}{内容与页面}
  \begin{columns}[T,onlytextwidth]
    \begin{column}{.46\textwidth}
      \textbf{大纲阶段}
      \begin{itemize}
        \item 组织章节与小节
        \item 补齐解释与证据
        \item 审阅并确认内容
      \end{itemize}
    \end{column}
    \begin{column}{.46\textwidth}
      \textbf{制作阶段}
      \begin{itemize}
        \item 划分页面与选择版式
        \item 分配正文与备注
        \item 检查并导出 PDF/PPTX
      \end{itemize}
    \end{column}
  \end{columns}
  \begin{notebox}[制作前提]
    详细大纲已确认；正文保留必要证据与限制。
  \end{notebox}
\end{frame}
\note{先对照左右两栏的大纲阶段与制作阶段，再指出下方“制作前提”说明框：它组织与主结构相关且必须可见的条件。条件用 notebox，操作建议用 tipbox，关键限制用 alertbox，定义和示例用对应环境。标题明确内容类别；不把每条正文都装箱，也不重复整页总结。}


\subsection{定义与示例}

\begin{frame}{定义与示例}
      \begin{definitionbox}[定义]
            \textbf{Modality}：信息的一种表示形式，如文本、图像、音频等。
            多模态学习关注如何联合建模这些异构来源。
      \end{definitionbox}

      \begin{examplebox}[示例]
            以 Tensor Fusion 为例，它通过外积显式保留单模态、双模态、
            三模态的交互，从而得到更丰富的联合表示。
      \end{examplebox}

\end{frame}
\note{比较定义与示例的语义。Tensor Fusion 文字沿用模板示例，正式论文讲解前应阅读并核实对应论文。}

\subsection{使用方法}

\begin{frame}[fragile]{提示框用法}
      \begin{verbatim}
\begin{tipbox}[自定义标题]
  需要强调的内容。
\end{tipbox}
      \end{verbatim}
      \begin{itemize}
            \item 将环境名替换为 \texttt{notebox}、\texttt{alertbox}、\texttt{examplebox} 或 \texttt{definitionbox}。
            \item 省略可选标题时，使用对应环境的默认标题。
      \end{itemize}
      \begin{tipbox}
            这是省略可选标题后的显示效果。
      \end{tipbox}
\end{frame}
\note{展示提示框环境的写法；省略标题会使用当前主题语言的默认名称。}

% !TeX root = ../example.tex
% 下面两个 section 专门覆盖 4、5 个 subsection 时的两栏排列。
\section{导航布局}
\sectioncoverpage[assets/tsinghua_door1]{本节用于检查四个 subsection 的两栏排列。} % 校园图片起始页
\note{本节检查四个小节在章节页上的交替双栏排列，以及页脚导航的联动。}
\subsection{第一主题}
\begin{frame}{四小节：第一主题}
  第 1 个 subsection 位于左列第一行。
\end{frame}
\note{第一小节对应章节页左列第一行。可点击导航验证目标。}
\subsection{第二主题}
\begin{frame}{四小节：第二主题}
  第 2 个 subsection 位于右列第一行。
\end{frame}
\note{第二小节对应章节页右列第一行，观察当前项高亮。}
\subsection{第三主题}
\begin{frame}{四小节：第三主题}
  第 3 个 subsection 位于左列第二行。
\end{frame}
\note{第三小节对应章节页左列第二行。}
\subsection{第四主题}
\begin{frame}{四小节：第四主题}
  第 4 个 subsection 位于右列第二行。
\end{frame}
\note{第四小节对应章节页右列第二行。接下来检查五小节与媒体页面。}

% !TeX root = ../example.tex
\section{媒体与收尾}
\sectioncoverpage[assets/sigs_night]{本节演示视频链接、参考文献与致谢。}
\subsection{视频链接}

\begin{frame}{视频链接示例}
      \centering
      \href{run:../assets/presentation_demo.mp4}{% 原创 CC0 动画；路径相对于 build/ 中的 PDF
            \XeTeXLinkBox{\fbox{\parbox[c][.30\textheight][c]{0.65\textwidth}{\centering 点击打开示例视频}}}%
      }
      \par\medskip
      {\small 原创动画：大纲 \(\rightarrow\) 幻灯片 \(\rightarrow\) 备注（6 秒，无音轨，CC0）。}
\end{frame}

\subsection{参考文献与致谢}

\begin{frame}[fragile]{参考文献与致谢用法}
      \begin{description}[汇总单页]
            \item[汇总单页] \texttt{\string\referencespage[shrink][\string\small]} 可指定字号并缩放到一页。
            \item[汇总续页] \texttt{\string\referencespage[break][\string\footnotesize]} 内容过多时自动续页。
            \item[双栏汇总] 命令末尾加 \texttt{[twocolumn]}；省略时默认为单栏。
            \item[致谢页面] \texttt{\string\backmatter} 会同步显示封面的标题、副标题、课程/课题组、作者、指导教师和日期；加入 \texttt{[notitle]} 可隐藏主标题。
      \end{description}
\end{frame}

\subsection{收尾说明}
\begin{frame}{收尾页选项}
  \begin{itemize}
    \item 参考文献页可以选择字号、续页和栏数。
    \item 致谢页可以保留或隐藏主标题。
  \end{itemize}
\end{frame}

\subsection{参考文献汇总}
\referencespage[break][\small][onecolumn] % 第三个参数改为 [twocolumn] 可使用双栏

\subsection{致谢}
\backmatter     % 不显示主标题时改为 \backmatter[notitle]

\end{document}
